\label{sec:theory}

Let $\mathbf v_k$ and $\mathbf w_k$ denote the right and left eigenvectors of $\mathbf T$ at non-unit eigenvalue $\lambda_k$, and let $\sigma_{|G|}^2$ denote the stationary variance of $|G_t|$ under the marginal mixture~\eqref{eq:mixture}. The closed-form bilinear identity for the absolute-return ACF of a stationary CHMM, $\E[|G_t|\,|G_{t+\tau}|] = \mathbf m^\top \mathrm{diag}(\bar{\boldsymbol\pi})\,\mathbf T^\tau\,\mathbf m$ with $\mathbf m_k = \E[|G_t| \mid s_t = k]$, and its eigendecomposition over the non-unit eigenvalues of $\mathbf T$ are folklore in the regime-switching literature (\citealp{hamilton1994time}, \S 22.2; \citealp{krolzig1997markov}, Ch.~3; \citealp{timmermann2000moments}). Under irreducibility and aperiodicity of $\mathbf T$ (Assumption~\ref{ass:irred}, with $\bar{\boldsymbol\pi}$ the unique stationary distribution; \citealp{levinperes2017markov}) and finite per-state second moments with non-trivial location and scale contrast (Assumption~\ref{ass:moments}), the spectral decomposition $\mathbf T^\tau = \mathbf 1\bar{\boldsymbol\pi}^\top + \sum_{k=2}^K \lambda_k^\tau\, \mathbf v_k \mathbf w_k^\top$ subtracts the marginal product to give the ACF
\begin{equation}
    \rho_{|G|}(\tau) \;=\; \sum_{k=2}^{K} w_k\, \lambda_k^{\tau},
    \qquad
    w_k = \big(\mathbf m^\top \mathrm{diag}(\bar{\boldsymbol\pi})\, \mathbf v_k\big)\big(\mathbf w_k^\top \mathbf m\big) / \sigma_{|G|}^2.
    \label{eq:acf_normalised}
\end{equation}
Equation~\eqref{eq:acf_normalised} is a real-coefficient sum of geometric decays at the non-unit eigenvalues of $\mathbf T$; complex-conjugate $\lambda_k$ pairs combine into damped oscillatory modes, and non-diagonalisable $\mathbf T$ adds polynomial-times-geometric prefactors via Jordan blocks. A self-contained derivation, the lag-zero behaviour, and the squared-return ACF are reported in the appendix, together with the formal Assumptions~\ref{ass:irred} and~\ref{ass:moments} and the cited propositions on identifiability, MLE consistency, and ECM monotonicity. Our contribution on this axis is exclusively empirical. We apply equation~\eqref{eq:acf_normalised} to recast the \citet{ryden1998stylized} low-$K$ failure on the SPY 2014 to 2024 instance as a rank statement on $\mathbf T - \mathbf 1\bar{\boldsymbol\pi}^\top$, and we demonstrate empirically that the implied $K - 1$ rank bound is non-binding at $K \ge 3$.

Equation~\eqref{eq:acf_normalised} reframes the Ryd\'{e}n low-$K$ failure as a rank statement on $\mathbf T - \mathbf 1\bar{\boldsymbol\pi}^\top$. At $K = 2$, equation~\eqref{eq:acf_normalised} reduces to the mono-exponential form $\rho_{|G|}(\tau) = \lambda_2^\tau$ with $\lambda_2 = T_{11} + T_{22} - 1$, and slow decay is matched by tuning $T_{11} + T_{22}$. The temporal axis is therefore satisfiable at $K = 2$; the binding axis is the two-component marginal, which cannot match the empirical kurtosis target. At $K \ge 3$ the matrix admits multiple non-unit eigenvalues and the marginal mixture has enough components to match the observed tail. Empirically, the absolute-return ACF-MAE was nearly flat across $K \in \{3, 6, \ldots, 21\}$ while distributional KS rose, and a direct effective-rank diagnostic showed a single non-unit eigenvalue carrying $93.6\%$ of the lag-1 ACF at $K = 18$ on SPY (Table~\ref{tab:spectral_rank}), so the $K$-rank upper bound is non-binding at $K \ge 3$ on this instance. Repeating the diagnostic on the 30-ticker cross-section gave a wider distribution with cross-ticker median dominant-mode share $0.76$ and minimum $0.326$ on NEM (Table~\ref{tab:spectral_xticker}); the rank-non-binding claim is supported at the cross-ticker median rather than as a universal property of equity-return data.
